%%%%%%%%%%%%%%%%%%%%%%%%%%%%%%%%%%%%%%%%%%%%%%%%%
%%%%%%%%%%%% chap: related work %%%%%%%%%%%%%%%%%
%%%%%%%%%%%%%%%%%%%%%%%%%%%%%%%%%%%%%%%%%%%%%%%%%

\chapter{Related Work}\label{chapter:related_work}

The research area relevant to this thesis sits at the intersection of retrieval-augmented
generation, generative engine optimization, and robustness against strategic manipulation.
Together, these directions show that external data is not just supplementary context for large
language models. It is a control surface that can improve factual grounding, alter visibility,
change ranking outcomes, and shape the final wording of generated answers.

\section{Retrieval-Augmented Generation as the Influence Surface}

The most general technical foundation for this thesis is retrieval-augmented generation
(RAG). Gao et al.~\cite{gao2024ragsurvey} describe RAG as a response to three well-known
limitations of large language models: hallucination, outdated knowledge, and non-transparent
reasoning. Instead of relying only on information stored in model weights, RAG systems
retrieve external documents and use them during generation. The survey emphasizes that this
combination improves accuracy and credibility, especially for knowledge-intensive tasks, while
also allowing continuous knowledge updates from external databases.

An important contribution of the survey is its organization of the field into Naive RAG,
Advanced RAG, and Modular RAG paradigms~\cite{gao2024ragsurvey}. This taxonomy is
useful for the present thesis because it makes clear that influence does not happen at a single
point. It can arise during retrieval, during augmentation of the prompt context, and during the
generation step itself. If the retrieval stage misses a relevant page, that page will not influence
the answer. If augmentation over-prioritizes one document, that document may dominate the
generated response. If generation compresses the evidence poorly, important facts may be
omitted even when they were retrieved correctly.

For this reason, RAG literature is directly relevant to the topic of external influence. It shows
why LLM outputs are sensitive to surrounding data and why document owners can affect the
model indirectly. However, the survey by Gao et al.~\cite{gao2024ragsurvey} is broad by
design. It explains how RAG systems work and how they are evaluated, but it does not focus
on website-side observability, content governance, or approval workflows. Those practical
concerns are central to this thesis.

\section{Generative Engine Optimization}

Aggarwal et al.~\cite{aggarwal2024geo} formalize the shift from traditional search engines to
``generative engines,'' systems that retrieve content and synthesize direct answers with large
language models. Their paper introduces Generative Engine Optimization (GEO) as a
creator-centric framework for improving how content appears in generated responses. This is a
key conceptual step because it reframes visibility as more than search ranking. In the GEO
view, a page can be evaluated by whether it is mentioned, how prominently it is surfaced, and
how much of its content is retained in the final answer.

The same work also introduces GEO-Bench, a benchmark of user queries and relevant web
sources designed for systematic evaluation~\cite{aggarwal2024geo}. The authors report that
their optimization strategies can increase visibility by up to 40\% and note that the effect varies
across domains. This finding matters for the present thesis in two ways. First, it demonstrates
that external data can measurably influence LLM-generated responses. Second, it suggests that
the influence is not uniform: different domains reward different document structures, phrasing
styles, and evidence patterns.

At the same time, GEO remains primarily a research framework for optimization and
evaluation. It does not focus on end-user operational concerns such as preserving locked facts,
requiring explicit approval before export, or presenting benchmark runs inside a product-style
dashboard. NewGEO builds on the same problem framing, but targets a practical workflow for
content teams rather than a benchmark-only environment.

\section{Strategic Manipulation of LLM Recommendations}

If GEO highlights the optimization opportunity, the work by Kumar and Lakkaraju
shows the risk side of the same phenomenon. Their paper demonstrates that an LLM's product
recommendations can be manipulated through a strategic text sequence (STS) added to a
product information page~\cite{kumar2024manipulating}. In their experiments with a
synthetic coffee-machine catalog, the injected sequence substantially increases the probability
that a target product becomes the top recommendation.

This paper is especially relevant because it makes the mechanism of external influence explicit.
The model is not retrained and the retrieval system is not replaced. Instead, the input content
itself is engineered so that the downstream LLM changes its ranking behavior. The authors
also emphasize robustness by optimizing the injected sequence against variations in the order in
which product descriptions appear in the prompt~\cite{kumar2024manipulating}. This
highlights an important insight: once an LLM-based ranking or recommendation system is
conditioned on external data, even small changes in phrasing can produce large competitive
effects.

From the perspective of this thesis, the contribution of Kumar and Lakkaraju is not that it
provides a deployment model to follow. On the contrary, it serves as evidence that any system
for GEO must distinguish between legitimate optimization and manipulative prompt shaping.
That distinction motivates the guarded design of NewGEO, where recommendations are paired
with constraint checks and a human approval step.

\section{Cooperative Optimization and Preference Learning}

More recent work pushes GEO beyond hand-written heuristics. Wu et al.~\cite{wu2025autogeo}
propose AutoGEO, a framework that learns generative-engine preference rules automatically
and uses them to rewrite documents. Their approach extracts preference rules from model
explanations, then applies those rules in two ways: as prompt context for a frontier-model
rewriter and as reward signals for a smaller trained model. This is important because it shifts
the research direction from manually selected optimization cues toward learned, engine-specific
rewriting behavior.

AutoGEO also introduces a more cooperative framing of optimization. Instead of evaluating
only visibility, the paper considers generative engine utility as well, asking whether rewriting
preserves response quality and reliability~\cite{wu2025autogeo}. The authors report an average
improvement of 35.99\% in GEO metrics while maintaining utility, and they observe that
different engines and domains prefer different rules. This directly supports one of the core
assumptions behind NewGEO: external influence should not be optimized blindly. It should be
measured against both visibility gains and semantic preservation.

Even so, AutoGEO remains focused on model-centric rewriting pipelines and benchmark
evaluation. It does not address the full operational lifecycle needed by a product team:
managing tracked pages, storing snapshots, defining query clusters, preserving locked facts,
approving recommendations, and exporting reviewed drafts. These are precisely the system-level
concerns that this thesis emphasizes.

\section{Positioning of the Thesis}

Taken together, the reviewed literature establishes four important points. First,
retrieval-augmented generation makes external knowledge a decisive part of the LLM response
process~\cite{gao2024ragsurvey}. Second, GEO research shows that content can be optimized to
gain visibility in generative answers~\cite{aggarwal2024geo}. Third, adversarial experiments
show that the same surface can be manipulated for competitive advantage~\cite{kumar2024manipulating}.
Fourth, newer work argues that optimization should be cooperative and utility-preserving, not
merely visibility-seeking~\cite{wu2025autogeo}.

The gap addressed by this thesis is practical and architectural. Existing work demonstrates that
external data can influence LLM outputs, but it leaves open how a real system should expose
that process to users in a controlled workflow. NewGEO answers this gap by providing an
inspectable platform where content can be imported, benchmarked, rewritten with supporting
context, validated against constraints, reviewed by a human, and exported only after approval.
In this sense, the thesis translates current GEO research into a concrete software architecture
for responsible experimentation.
